% Guía rápida TeleCon4ESP32 — solo la primera sesión Bluetooth.
% Compilar: pdflatex -output-directory=../build telecon_fast_guide.tex
% PDF entregado (asset Android Codes): ../build/fast_guide_esp.pdf
%   (copiar y quitar telecon_fast_guide.pdf — ver build_guides.sh)
\documentclass[12pt,a4paper]{article}

\usepackage[margin=2.4cm]{geometry}
\usepackage[T1]{fontenc}
\usepackage[utf8]{inputenc}
\usepackage{lmodern}
\usepackage[spanish]{babel}
\usepackage{booktabs}
\usepackage{hyperref}
\usepackage{xcolor}
\usepackage{colortbl}
\usepackage{array}
\usepackage{enumitem}
\usepackage{needspace}
\usepackage{etoolbox}
\usepackage{tikz}
\usetikzlibrary{shadows}
\usepackage{graphicx}
\usepackage{float}
\usepackage{caption}
\captionsetup{justification=centering,singlelinecheck=false}

\preto\section{\Needspace{8\baselineskip}}

\hypersetup{
  colorlinks=true,
  linkcolor=black,
  urlcolor=blue,
  pdftitle={Guía rápida TeleCon4ESP32},
  pdfauthor={TeleCon}
}

\definecolor{ArduinoKeyword}{HTML}{00979C}
\definecolor{ArduinoFunc}{HTML}{D35400}
\definecolor{ArduinoComment}{HTML}{5D6D6E}
\definecolor{ArduinoPreproc}{HTML}{9B59B6}
\definecolor{TblHead}{HTML}{0F6F73}
\definecolor{TblChan}{HTML}{D5F5E3}
\definecolor{TblSw}{HTML}{FCF3CF}
\definecolor{TblSum}{HTML}{FAD7A0}
\definecolor{TblTelem}{HTML}{D4E6F1}
\definecolor{TblSync}{HTML}{E8DAEF}
\definecolor{TblId}{HTML}{F5CBA7}

\newcommand{\thc}[1]{\textcolor{white}{\bfseries #1}}
\newcommand{\chip}[2]{\colorbox{#1}{\strut\,\scriptsize #2\,}}
\newcommand{\ui}[1]{\textsf{\textbf{#1}}}

\newcommand{\shot}[4]{%
  \par\addvspace{0.25em}%
  \IfFileExists{../../figures/#2}{%
    \begin{figure}[H]
      \centering
      \begin{tikzpicture}
        \node[inner sep=0] (img) {%
          \phantom{\includegraphics[width=0.68\textwidth,height=0.48\textheight,keepaspectratio]{../../figures/#2}}%
        };
        \fill[black, fill opacity=0.12, rounded corners=13pt]
          ([xshift=-5.5pt+1.8mm,yshift=-5.5pt-2.3mm]img.south west) rectangle
          ([xshift=5.5pt+1.8mm,yshift=5.5pt-2.3mm]img.north east);
        \fill[black!92, rounded corners=13pt]
          ([xshift=-5.5pt,yshift=-5.5pt]img.south west) rectangle
          ([xshift=5.5pt,yshift=5.5pt]img.north east);
        \begin{scope}
          \clip[rounded corners=10pt]
            (img.south west) rectangle (img.north east);
          \node[inner sep=0] at (img.center) {%
            \includegraphics[width=0.68\textwidth,height=0.48\textheight,keepaspectratio]{../../figures/#2}%
          };
        \end{scope}
        \draw[black!30, line width=0.65pt, rounded corners=13pt]
          ([xshift=-5.5pt,yshift=-5.5pt]img.south west) rectangle
          ([xshift=5.5pt,yshift=5.5pt]img.north east);
        \draw[TblHead, line width=1.25pt, line cap=round]
          ([xshift=-5.5pt+11pt,yshift=-5.5pt+3.4pt]img.south west) --
          ([xshift=5.5pt-11pt,yshift=-5.5pt+3.4pt]img.south east);
      \end{tikzpicture}%
      \caption{\textbf{#1.} #4\\[0.35em]
        {\sffamily\scriptsize\color{ArduinoComment}#3
          \enspace\textbullet\enspace\texttt{\detokenize{#2}}}}
    \end{figure}%
  }{%
    \noindent
    \begin{tikzpicture}
      \node[draw=TblHead, fill=TblSw, line width=1.1pt, rounded corners=6pt,
            inner sep=9pt, text width=\dimexpr\linewidth-20pt\relax, align=left] {
        {\sffamily\bfseries\color{TblHead} Captura #1
          \hfill{\normalfont\ttfamily\scriptsize \detokenize{#2}}}\\[0.3em]
        {\sffamily\scriptsize\color{ArduinoFunc} #3}\\[0.35em]
        {\sffamily\small #4}\\[0.3em]
        {\sffamily\scriptsize\color{ArduinoComment}
          Guardar como \texttt{General\_Documentation/figures/\detokenize{#2}}.}
      };
    \end{tikzpicture}%
  }%
  \par\addvspace{0.2em}%
}

\title{Guía rápida TeleCon4ESP32\\[0.3em]
\large Primera sesión: Predeterminado + Classic Simple}
\author{Teléfono y ESP32 DevKit --- unos diez minutos\\[0.2em]
\small Recorrido completo: PDF de la guía de usuario}
\date{Versi\'on 1.0 \quad 2026}

\begin{document}
\maketitle

\begin{abstract}
Mueve un stick en \ui{Panel de control} con Bluetooth Classic Simple.
Sin Pro, sin monedas, sin cámara. La guía de usuario cubre Vehículo RC,
Wi-Fi y fallos.
\end{abstract}

{\sffamily\scriptsize
\chip{TblChan}{Usuario predeterminado}
\chip{TblSw}{Classic Simple}
\chip{TblSum}{Panel de control --- gratis}}

\section{Qué necesitas}
\begin{itemize}[leftmargin=1.4cm]
  \item Teléfono Android con TeleCon4ESP32.
  \item ESP32 DevKit (sin cámara), cable USB, Arduino IDE.
\end{itemize}

\section{Cargar la placa}
Abre la carpeta que contiene \texttt{TeleCon\_Classic\_Simple.ino}.
Herramientas $\rightarrow$ Placa $\rightarrow$ \textbf{ESP32 Dev Module}.
Elige el puerto USB. Sube el programa. Monitor Serie \textbf{115200}.
Pulsa RESET. Espera el banner de listo. La placa anuncia
\texttt{ESP32-TC-RC-BT-Simple}.

\shot{S01}{s01_serial.png}{Ordenador, Monitor Serie de Arduino IDE}
{Serial a 115200 tras RESET. Banner de listo visible.}

\section{Ajustes del teléfono}
Concede Dispositivos cercanos / Bluetooth cuando Android lo pida.
Abre \ui{Panel de control} $\rightarrow$ engranaje:

\begin{center}
\renewcommand{\arraystretch}{1.22}
\begin{tabular}{>{\raggedright\arraybackslash}p{4.2cm}
                >{\raggedright\arraybackslash}p{8.0cm}}
\rowcolor{TblHead}
\thc{Ajuste} & \thc{Valor} \\
\rowcolor{TblChan}
Tipo de usuario & Predeterminado \\
Placa & DevKit \\
Conexión & Bluetooth, Classic + Simple (líneas de texto) \\
\end{tabular}
\end{center}

Sal de ajustes. Toca Bluetooth en la barra del módulo. Escanea.
Toca el nombre anunciado
\texttt{ESP32-TC-RC-BT-Simple}
(PIN \texttt{1234} si lo pide). El estado debe mostrar \ui{Conectado}.

\shot{S02}{s02_settings.png}{Vertical, tema claro}
{Ajustes de Panel de control: Predeterminado + Classic + Simple
(líneas de texto). Interfaz en español.}

\shot{S03}{s03_scan.png}{Vertical, escaneo Bluetooth}
{ESP32-TC-RC-BT-Simple en Dispositivos escaneados. Interfaz en español.}

\section{Mover un stick}
Abre \ui{Panel de control} (horizontal). Arrastra el stick izquierdo.
Hasta que la placa responda OK, el movimiento en pantalla no es una
sesión en vivo --- aunque Android diga que Bluetooth está conectado.
En el banco, el firmware oficial hace eco o actualiza gauges sin motores.

Trata los sticks como vivos cuando ponga \ui{Conectado}. Los sketches
oficiales hacen fail-safe al caer el enlace (sticks a 0).

\section{Si falla}
Alimentación + RESET. Dispositivos cercanos permitidos. El modo coincide
con este sketch (no BLE, no Binary). Olvida un vínculo viejo y escanea
otra vez. Tabla completa: PDF de la guía de usuario.

\section{Después}
\begin{itemize}[leftmargin=1.4cm]
  \item PDF de la guía de usuario --- Inicio, Vehículo RC, roles de cámara.
  \item \texttt{README.md} del sketch --- pines, motores, trim.
  \item No pases a Avanzado hasta que este camino funcione.
\end{itemize}

\end{document}
